\subsection*{最初に必要な用語}
\addcontentsline{toc}{subsection}{最初に必要な用語}

\par\smallskip\noindent\refstepcounter{table}\label{tab:ch15-01}表 \thetable: 最初に必要な用語の整理\par\smallskip
表 \thetable は、最初に必要な用語の整理を比較・確認しやすい形で整理したものである。\par\smallskip
\begin{longtable}{L{0.27\linewidth} L{0.68\linewidth}}
\toprule
用語 & 初学者向けの説明 \\
\midrule
ソフトウェア工学経済 & 技術案を費用、価値、リスク、時間、組織目標の観点から比較し、よりよい選択を行うための知識である。 \\
提案 & 実行するかどうかを判断する一つの選択肢である。例は「新規開発する」「既存製品を購入する」「レガシーシステムを置き換える」である。 \\
キャッシュフロー & 提案を実行した結果、いつ、いくらの資金が出入りするかを表す流れである。 \\
時間価値 & 同じ金額でも、今日の金額と将来の金額では価値が異なるという考え方である。 \\
等価性 & 異なる時点の金額を、同じ時点の価値に換算して比較できる状態である。 \\
比較基準 & 提案を同じ土俵で比べるための基準である。現在価値、将来価値、内部収益率などがある。 \\
\term{無形資産}{Intangible Asset} & 物理的な形はないが価値を持つ知識、文化、手順、経験、ノウハウ、仕様、ツールなどである。 \\
見積り & まだ正確には分からない量を、根拠に基づいて予測することである。 \\
総所有費用 & 取得費だけでなく、導入、運用、保守、廃止までを含めた総費用である。 \\
\bottomrule
\end{longtable}
